\uuid{ymJ0}
\exo7id{691}
\titre{exo7 691}
\auteur{bodin}
\organisation{exo7}
\datecreate{1998-09-01}
\isIndication{false}
\isCorrection{false}
\chapitre{Continuité, limite et étude de fonctions réelles}
\sousChapitre{Etude de fonctions}
\cpge{mpsi-mp2i.limites-continuite.c, pcsi.limites-continuite.c, ptsi.limites-continuite.c, mpsi-mp2i.fonctions-valeurs-reelles-complexes.b, pcsi.techniques-calcul-fonctions.b, ptsi.techniques-calcul-fonctions.b, mpsi-mp2i.fonctions-valeurs-reelles-complexes.a, pcsi.techniques-calcul-fonctions.a, ptsi.techniques-calcul-fonctions.a, bcpst1.limites-continuite.d, bcpst1.derivation.a, bcpst1.fonctions-reelles-usuelles.a}
\module{Analyse}
\niveau{L1}
\difficulte{}

\contenu{
\begin{enumerate}
    \item \question{Soit la fonction $f:\lbrack -1,+\infty\lbrack\to\R$, d\'efinie par
$f(x)=\frac{1}{\sqrt{x^2+2x+2}}$. Montrer que $f$ admet une r\'eciproque
 que l'on explicitera.}
    \item \question{Trouver un intervalle de $\R$ sur lequel la fonction $g(x)=\tan(x^3)$
admette une fonction r\'eciproque (on pr\'ecisera alors le domaine de
d\'efinition de cette r\'eciproque et son image).}
\end{enumerate}
}
